% © 2026 Ing. David Jaroš — CC BY-NC-ND 4.0
%
% This work is licensed under a Creative Commons Attribution-NonCommercial-NoDerivatives
% 4.0 International License (CC BY-NC-ND 4.0).
%
% License History: Earlier drafts (up to v0.3) were released under CC BY 4.0.
% From v0.4 onward, all material is released under CC BY-NC-ND 4.0 to protect
% the integrity of the theoretical work during ongoing academic development.
%
% See LICENSE.md for full license text.
%
% Track: T3_ALPHA — Fine Structure Constant
% Purpose: Path A of the Gap G137-B attack — Modular bootstrap on the partition
%          function Z(τ) = ϑ₃³(τ) to derive B_phenom from crossing symmetry.
%          This extends the earlier modular_bootstrap_k1_attempt.tex by applying
%          crossing symmetry to the 4-point function on T².
% Predecessor: research_tracks/T3_ALPHA/modular_bootstrap_k1_attempt.tex
% Status: RESEARCH DRAFT — path A attempt; go/no-go gate 2026-05-27
% Date: 2026-05-10

\documentclass[12pt,a4paper]{article}
\usepackage[a4paper,margin=2.5cm]{geometry}
\usepackage{amsmath,amssymb,amsthm}
\usepackage{mathtools}
\usepackage{hyperref}
\usepackage{tcolorbox}
\tcbuselibrary{skins,breakable}
\usepackage{booktabs}
\usepackage{xcolor}

\hypersetup{
  colorlinks=true,
  linkcolor=blue!60!black,
  citecolor=green!50!black
}

\newtheorem{theorem}{Theorem}[section]
\newtheorem{lemma}[theorem]{Lemma}
\newtheorem{proposition}[theorem]{Proposition}
\newtheorem{corollary}[theorem]{Corollary}
\theoremstyle{definition}
\newtheorem{definition}[theorem]{Definition}
\theoremstyle{remark}
\newtheorem{remark}[theorem]{Remark}

\newcommand{\PROVED}{\textbf{[Proved]}}
\newcommand{\OPEN}{\textbf{[Open]}}
\newcommand{\GAP}[1]{\textbf{[Gap~#1]}}
\newcommand{\MC}{\textbf{[MC]}}
\newcommand{\Lzero}{\textbf{[L0]}}
\newcommand{\Lone}{\textbf{[L1]}}

\newcommand{\CC}{\mathbb{C}}
\newcommand{\HH}{\mathbb{H}}
\newcommand{\ZZ}{\mathbb{Z}}
\newcommand{\RR}{\mathbb{R}}
\newcommand{\SL}{\mathrm{SL}(2,\mathbb{Z})}
\newcommand{\vth}{\vartheta}
\newcommand{\calH}{\mathcal{H}}

\title{\textbf{Path A: Modular Bootstrap on $\hat{Z}(\tau) = \vth_3^3(\tau)$}\\[6pt]
  \large Crossing Symmetry on $T^2$ and the Kac-Moody Level $k_{\mathrm{KM}} = 1$\\[4pt]
  \normalsize T3\_ALPHA — Gap G137-B Attack, Path A}
\author{Ing.\ David Jaroš}
\date{2026-05-10 — Time-box: 4 weeks (gate: 2026-05-27)}

\begin{document}
\maketitle

\begin{tcolorbox}[enhanced,colback=orange!5!white,colframe=orange!60!black,
  arc=3pt,boxrule=1.2pt,title={\textbf{Hard Rules (Non-Negotiable, from ALPHA\_BREAKTHROUGH\_REPORT.md)}}]
\textbf{R1}: No number may be fitted to reproduce $\alpha$ or 137.\\
\textbf{R2}: Every constant must originate from an independent UBT sector.\\
\textbf{R3}: No experimental input ($\alpha_{\exp}$, 137, $B_{\mathrm{required}}$) may appear.\\
\textbf{R4}: Every step is classified: CLEAN [L0/L1], MC, or SE.\\
\textbf{R5}: No CIRC (circular) input may appear in a first-principles claim.
\end{tcolorbox}

\tableofcontents

% ==============================================================================
\section{Objective}
\label{sec:obj}
% ==============================================================================

The single blocking step to the integer-137 result in UBT is Gap~G137-B:
derive $B = B_{\mathrm{phenom}} \approx 46.298$ from the UBT action $S[\Theta]$
without any experimental input.

The predecessor document \texttt{modular\_bootstrap\_k1\_attempt.tex} identified
the Kac-Moody level $k_{\mathrm{KM}} = 1$ as a key sub-gap (Gap G3-k): if
$k_{\mathrm{KM}} = 1$ is forced by consistency, then
$B_{\mathrm{base}} = N_{\mathrm{eff}}^{3/2} \approx 41.57$, and the remaining
gap to $B_{\mathrm{phenom}}$ can be closed by the elliptic-point correction
$\varepsilon$ from $\Gamma_0(137)$.

This document attempts \textbf{Path A}: impose crossing symmetry on the 4-point
correlation function of $\Theta$ on the torus $T^2 = \CC/(\ZZ + \tau\ZZ)$ and
determine whether it forces $k_{\mathrm{KM}} = 1$.

% ==============================================================================
\section{Setup and Pre-Proved Inputs}
\label{sec:setup}
% ==============================================================================

The following results are used from earlier proofs; they must not be re-derived here.

\begin{tcolorbox}[colback=green!5!white,colframe=green!50!black,
  title={Pre-Proved Inputs (must not be re-derived)}]
\begin{center}
\begin{tabular}{lll}
\toprule
\textbf{Result} & \textbf{Level} & \textbf{Source} \\
\midrule
$N_{\mathrm{eff}} = 12$ & \Lzero & \texttt{canonical/n\_eff/step3\_N\_eff\_result.tex} \\
$V_{\mathrm{eff}}(n) = n^2 - Bn\ln n$ & \Lone\ (given $B$) & \texttt{canonical/alpha/alpha\_best\_route.tex} \\
$n^*(B_{\mathrm{phenom}}) = 137$ & \Lone\ (given $B$) & same \\
$B_0 = 8\pi$ (one-loop) & \Lone & \texttt{canonical/n\_eff/step2\_vacuum\_polarization.tex} \\
$\hat{Z}(\tau) = \vth_3^3(\tau)$, modular weight $3/2$ & \Lone & \texttt{ALPHA\_PROGRESS\_REPORT.md §3.4} \\
$\mu(\Gamma_0(137))/3 = 46.000$ & [STD] & \texttt{reports/gamma0\_137\_invariants.md} \\
\bottomrule
\end{tabular}
\end{center}
\end{tcolorbox}

\subsection{Torus geometry}

The imaginary-time circle $S^1_\psi$ of radius $R_\psi$ is compactified.
The UBT Euclidean field theory lives on
\[
  T^2 \;=\; \CC\,/\,(\ZZ + \tau\ZZ),
  \qquad
  \tau \;=\; \frac{iR_\psi}{R_t} \;\in\; \mathbb{H}_+,
\]
where $\mathbb{H}_+ = \{\tau \in \CC : \mathrm{Im}\,\tau > 0\}$.

\subsection{Partition function and modular content}

The UBT partition function is established to be \Lone:
\[
  \hat{Z}(\tau) \;=\; \vth_3^3(\tau)
  \;=\;
  \Bigl(\sum_{n\in\ZZ} e^{i\pi\tau n^2}\Bigr)^3.
\]
Under modular $S$-transformation $\tau \to -1/\tau$:
\[
  \hat{Z}(-1/\tau)
  = \vth_3^3(-1/\tau)
  = (-i\tau)^{3/2}\,\vth_3^3(\tau)
  = (-i\tau)^{3/2}\,\hat{Z}(\tau).
\]
The modular weight is $h = 3/2$, corresponding to $N_{\mathrm{eff}}/8 = 12/8 = 3/2$
real scalars (each contributing weight $1/2$).

% ==============================================================================
\section{The 4-Point Correlation Function on $T^2$}
\label{sec:4pt}
% ==============================================================================

\subsection{Winding-mode vertex operators}

Winding modes at level $n$ correspond to vertex operators
\[
  V_n(\sigma) \;=\; e^{i n X_\psi(\sigma)},
\]
where $X_\psi(\sigma)$ is the imaginary-time boson at world-sheet position
$\sigma \in T^2$, and $n \in \ZZ$ is the winding number.  The conformal
weight of $V_n$ in the $\hat{u}(1)^{N_{\mathrm{eff}}}$ current algebra at
Kac-Moody level $k$ is
\begin{equation}
\label{eq:confwt}
  h_n \;=\; \frac{n^2}{4k}\,\cdot\,\frac{N_{\mathrm{eff}}}{N_{\mathrm{eff}}} = \frac{n^2}{4k}.
\end{equation}
For $k = 1$: $h_n = n^2/4$.  The corresponding $V_{\mathrm{eff}}$ mass term
$\sim n^2$ is consistent with $h_n \sim n^2$ at $k = 1$.

\subsection{4-point function on $T^2$}

The 4-point function of winding-mode operators at positions
$\sigma_1, \sigma_2, \sigma_3, \sigma_4 \in T^2$ is
\[
  G_4(\sigma_1,\sigma_2,\sigma_3,\sigma_4)
  \;=\;
  \bigl\langle V_{n_1}(\sigma_1)\,V_{n_2}(\sigma_2)\,
              V_{n_3}(\sigma_3)\,V_{n_4}(\sigma_4) \bigr\rangle_{T^2}.
\]
Charge conservation on $T^2$ requires $n_1 + n_2 + n_3 + n_4 = 0$.

The simplest non-trivial case is $n_1 = n_2 = n$, $n_3 = n_4 = -n$
(elastic winding-mode scattering):
\[
  G_4^{(n)}
  \;=\;
  \bigl\langle V_n(\sigma_1)\,V_n(\sigma_2)\,
              V_{-n}(\sigma_3)\,V_{-n}(\sigma_4) \bigr\rangle_{T^2}.
\]

% ==============================================================================
\section{Crossing Symmetry Constraint}
\label{sec:crossing}
% ==============================================================================

\subsection{Statement of crossing symmetry}

On $T^2$, modular invariance of the correlation function under
$\tau \to -1/\tau$ (S-transform) and $\tau \to \tau + 1$ (T-transform)
constitutes the analogue of crossing symmetry in the torus setting.

For the 4-point function $G_4^{(n)}$, crossing symmetry in the $s$- and
$t$-channel decomposition requires:

\begin{equation}
\label{eq:crossing}
  G_4^{(n)}(\tau) \;=\; G_4^{(n)}(-1/\tau).
\end{equation}

This is a strong constraint on the operator content of the theory.

\subsection{Virasoro block decomposition}

Inserting a complete set of intermediate states in the $s$-channel (winding
sector $m$):
\[
  G_4^{(n)}(\tau)
  \;=\;
  \sum_m C_{nm\bar{m}}\,C_{\bar{m}nm}\,
  \calF\Bigl(h_m; \frac{c}{24}; \tau\Bigr),
\]
where $\calF(h_m; c/24; \tau)$ is the Virasoro conformal block at
intermediate weight $h_m$ and central charge $c = N_{\mathrm{eff}} = 12$.

\subsection{The $k_{\mathrm{KM}} = 1$ constraint}

\begin{proposition}[Kac-Moody level constraint from crossing, \MC]
\label{prop:k1}
If the partition function $\hat{Z}(\tau) = \vth_3^3(\tau)$ is the unique
modular-covariant partition function of the UBT torus CFT compatible with
the central charge $c = N_{\mathrm{eff}} = 12$ and the winding-mode spectrum
$h_n \sim n^2$, then the crossing symmetry constraint~\eqref{eq:crossing}
forces $k_{\mathrm{KM}} = 1$.
\end{proposition}

\begin{tcolorbox}[colback=yellow!5!white,colframe=orange!60!black,
  title={Status of Proposition~\ref{prop:k1}: OPEN [MC]}]
The proposition is \textbf{not yet proved}.  The gap is:

\begin{enumerate}
  \item The Virasoro block decomposition of $\vth_3^3(\tau)$ has not been
        computed explicitly.  This requires knowledge of the OPE coefficients
        $C_{nm\bar{m}}$ in the winding sector.

  \item The uniqueness of $\hat{Z}(\tau) = \vth_3^3(\tau)$ among
        modular-covariant functions of weight $3/2$ is established (weight $3/2$
        forces $\vth_3^3$ among holomorphic weight-$3/2$ modular forms), but
        whether this uniqueness constrains $k_{\mathrm{KM}}$ to be $1$
        specifically has not been shown.

  \item The crossing equation~\eqref{eq:crossing} on $T^2$ imposes an
        infinite tower of consistency conditions (one per intermediate
        operator); checking these systematically requires numerical
        bootstrap methods (e.g., semi-definite programming as in the
        Rattazzi-Rychkov-Tonni-Vichi approach).
\end{enumerate}

\textbf{Classification}: \MC\ — the logic is sound in outline but the
detailed consistency check is not complete.  No experimental input is used.
\end{tcolorbox}

% ==============================================================================
\section{Consequence: $B_{\mathrm{base}} = N_{\mathrm{eff}}^{3/2}$ if $k = 1$}
\label{sec:bbase}
% ==============================================================================

\begin{lemma}[$B_{\mathrm{base}}$ from $k_{\mathrm{KM}} = 1$, \MC\ conditional on
  Proposition~\ref{prop:k1}]
\label{lem:bbase}
Assuming $k_{\mathrm{KM}} = 1$ for the $\hat{u}(1)^{N_{\mathrm{eff}}}$ current
algebra of the UBT torus, the effective coupling in $V_{\mathrm{eff}}$ is
\[
  B_{\mathrm{base}}
  \;=\;
  \frac{N_{\mathrm{eff}}^{3/2}}{(k_{\mathrm{KM}} + h^\vee)^{1/2}}
  \;\stackrel{k=1,\,h^\vee=0}{=}\;
  N_{\mathrm{eff}}^{3/2}
  \;=\;
  12^{3/2}
  \;\approx\;
  41.569.
\]
\end{lemma}

\begin{proof}[Proof sketch]
For a $\hat{u}(1)$ current algebra ($h^\vee = 0$) at Kac-Moody level $k$,
the conformal weight of the winding-mode state at level $n$ is
$h_n = n^2/(2k)$.  The $n\ln n$ coefficient in $V_{\mathrm{eff}}$ arises
from the functional determinant of the heat kernel of the Laplacian on
$S^1_\psi$, which is proportional to $N_{\mathrm{eff}} \cdot \sqrt{h_n/n^2}
= N_{\mathrm{eff}}/\sqrt{2k}$.
Setting $k = 1$ and using $N_{\mathrm{eff}} = 12$ gives
$B_{\mathrm{base}} = N_{\mathrm{eff}}^{3/2} = 41.57$.
(See \texttt{canonical/alpha/neff\_32\_alpha\_route.tex} for the full argument.)
\end{proof}

% ==============================================================================
\section{Elliptic-Point Correction to Close the Remaining Gap}
\label{sec:elliptic}
% ==============================================================================

Even if Proposition~\ref{prop:k1} is confirmed, there remains a gap:
\[
  B_{\mathrm{base}} = 41.569 \;\neq\; B_{\mathrm{phenom}} = 46.298.
\]
The ratio is $B_{\mathrm{phenom}}/B_{\mathrm{base}} \approx 1.114$.

From \texttt{reports/gamma0\_137\_invariants.md}:
\[
  \frac{\mu(\Gamma_0(137))}{3} \;=\; 46.000,
  \qquad
  \frac{\mu(\Gamma_0(137))}{3} + \frac{\nu_2}{4} \;=\; 46.500,
\]
where $\nu_2 = 2$ is the number of order-2 elliptic points of $\Gamma_0(137)$.

\begin{remark}[Residual gap]
The modular bootstrap sets $B_{\mathrm{base}} \approx 41.6$ (if $k=1$ is
confirmed), while $B_{\mathrm{phenom}} \approx 46.3$.  The remaining gap is
$(46.3 - 41.6)/41.6 \approx 11\%$.  The elliptic-point correction from
$\Gamma_0(137)$ contributes $\nu_2/4 = 0.5$, closing the gap only from $0.64\%$
to $0.44\%$ \emph{at} $B = (p+1)/3$, not from $10\%$ at $B = N_{\mathrm{eff}}^{3/2}$.

This means that confirming $k_{\mathrm{KM}} = 1$ via crossing symmetry gives
$B_{\mathrm{base}} \approx 41.6$, not $46.3$.  A further mechanism is required
to bridge from $41.6$ to $46.3$.  The elliptic correction alone is insufficient.
\end{remark}

% ==============================================================================
\section{Outcome Assessment}
\label{sec:outcome}
% ==============================================================================

\begin{tcolorbox}[enhanced,colback=red!4!white,colframe=red!60!black,
  arc=3pt,boxrule=1.2pt,title={\textbf{Path A Outcome Assessment (2026-05-10)}}]

\begin{center}
\begin{tabular}{ll}
\toprule
\textbf{Component} & \textbf{Status} \\
\midrule
Partition function $\hat{Z} = \vth_3^3(\tau)$ established & \PROVED\ \Lone \\
$c = N_{\mathrm{eff}} = 12$ consistent with $k=1$ & consistent [not forced] \\
Crossing symmetry constraint \eqref{eq:crossing} stated & Yes (this document) \\
Virasoro block decomposition completed & \OPEN \\
$k_{\mathrm{KM}} = 1$ forced by crossing (Prop.~\ref{prop:k1}) & \OPEN\ [MC] \\
$B_{\mathrm{base}} = N_{\mathrm{eff}}^{3/2} \approx 41.6$ (if $k=1$) & \MC\ conditional \\
Bridge from 41.6 to 46.3 & \OPEN\ [no mechanism found] \\
\midrule
\textbf{Path A resolves Gap G137-B} & \textbf{NO} (as of 2026-05-10) \\
\bottomrule
\end{tabular}
\end{center}

\smallskip
\textbf{Conclusion}: Path A clarifies the structure of the crossing constraint
and identifies $k_{\mathrm{KM}} = 1$ as a necessary (but not sufficient)
sub-goal.  Even if $k = 1$ is confirmed, an additional mechanism is required
to bridge $B_{\mathrm{base}} \approx 41.6$ to $B_{\mathrm{phenom}} \approx 46.3$.
No experimental input was used.

\textbf{Recommendation}: Continue numerical bootstrap methods (semi-definite
programming) to check the crossing constraint systematically.  Simultaneously
pursue Path B ($\zeta$-function regularisation).
\end{tcolorbox}

% ==============================================================================
\section{Next Steps for Path A}
\label{sec:nextsteps}
% ==============================================================================

To complete the crossing symmetry argument:

\begin{enumerate}
  \item Compute the OPE coefficients $C_{n_1 n_2 m}$ in the winding-mode
        sector from the $\vth_3^3(\tau)$ structure.

  \item Expand $G_4^{(n)}$ in Virasoro conformal blocks and impose the
        $s \leftrightarrow t$ crossing equation using the numerical bootstrap
        approach of~\cite{rattazzi2008}.

  \item Check whether the system is consistent only for $k_{\mathrm{KM}} = 1$
        (ruling out $k = 2, 3, \ldots$) within the constraints forced by
        $c = 12$ and the winding spectrum.

  \item If $k = 1$ is forced and a correction mechanism from $41.6$ to $46.3$
        is found, promote the result to \Lone\ and open a PR to
        \texttt{canonical/alpha/}.
\end{enumerate}

\begin{thebibliography}{9}

\bibitem{rattazzi2008}
R.~Rattazzi, V.~S.~Rychkov, E.~Tonni, and A.~Vichi,
\emph{Bounding scalar operator dimensions in 4D CFT},
JHEP \textbf{12}, 031 (2008).

\bibitem{diamond_shurman}
F.~Diamond and J.~Shurman,
\emph{A First Course in Modular Forms}, Springer, 2005.

\bibitem{ginsparg1988}
P.~Ginsparg, \emph{Applied Conformal Field Theory},
Les Houches Lectures 1988, hep-th/9108028.

\end{thebibliography}

\end{document}
